\documentclass[luatex,fontsize=10pt,paper=a5]{jlreq}
\usepackage{gckanbun,lltjext}

\begin{document}

\section*{gckanbun sample}

このサンプルは LuaLaTeX でそのまま試せる最小例です。

\subsection*{基本の訓点}

\begin{GCKEnv}{2\zw}
  \振り{学}{まな}\送り{ビテ}\返り{レ}時\送り{ニ}習\送り{フ}
\end{GCKEnv}

\subsection*{\textbackslash KanHyphen}

\begin{GCKEnv}{2\zw}
  \振り{所}{ゆ}\返り[intrusion=post]{二}\振り{\KanHyphen}{ゑ}\振り{以}{ん}
\end{GCKEnv}

\subsection*{グループルビ}

\begin{GCKEnv}{2\zw}
  \グ振り{所以}{ゆえん}\送り{ナリ}

  \bigskip

  \グ振り{所\返り[intrusion=post]{一}\KanHyphen 以}{ゆえん}\送り{ノミ}

  \bigskip

  \グ振り{所以}{ゆえん}[かくか]\送り{ナリ}[シム]
\end{GCKEnv}

\subsection*{局所縦組}

横組の文中でも、次のように一部分だけ縦組にできます。

\parbox<t>{8\zw}{\GCKTateOn
  \begin{GCKEnv}{2\zw}
    \振り{雖}{いへど}\送り{モ}\\
    \振り{所}{ゆ}\返り[intrusion=post]{二}\振り{\KanHyphen}{ゑ}\振り{以}{ん}\\
    \グ振り{以心}{いしん}\グ振り{伝心}{でんしん}
  \end{GCKEnv}
\par}

\end{document}
